\documentclass[aspectratio=169, 12pt]{beamer}

% ── Tema ────────────────────────────────────────────────────────────────────
\usetheme{Madrid}
\usecolortheme{default}
\setbeamercolor{palette primary}{bg=black!85, fg=white}
\setbeamercolor{palette secondary}{bg=black!70, fg=white}
\setbeamercolor{palette tertiary}{bg=black!60, fg=white}
\setbeamercolor{palette quaternary}{bg=black!50, fg=white}
\setbeamercolor{structure}{fg=black!80}
\setbeamercolor{block title}{bg=black!80, fg=white}
\setbeamercolor{block body}{bg=black!5}
\setbeamercolor{alerted text}{fg=orange!80!black}
\setbeamerfont{block title}{series=\bfseries}
\setbeamerfont{title}{size=\Large, series=\bfseries}
\setbeamerfont{frametitle}{series=\bfseries}
\setbeamertemplate{navigation symbols}{}
\setbeamertemplate{footline}[frame number]

% ── Paquetes ─────────────────────────────────────────────────────────────────
\usepackage[spanish]{babel}
\usepackage[utf8]{inputenc}
\usepackage{amsmath, amssymb}
\usepackage{booktabs}
\usepackage{tikz}
\usetikzlibrary{arrows.meta, positioning, shapes.geometric, fit, backgrounds}
\usepackage{xcolor}
\usepackage{graphicx}
\usepackage{fontawesome5}

% ── Colores custom ───────────────────────────────────────────────────────────
\definecolor{phaseA}{RGB}{52,  120, 180}
\definecolor{phaseB}{RGB}{200,  80,  40}
\definecolor{phaseC}{RGB}{  0, 150,  80}
\definecolor{phaseD}{RGB}{150,  60, 200}
\definecolor{darkgray}{RGB}{50,50,50}

% ── Metadata ─────────────────────────────────────────────────────────────────
\title{\textbf{DreamerV4}\\[4pt]
       \large Aprendizaje de Modelos de Mundo para Control de Largo Plazo}
\subtitle{Hansen et al., 2025 — \texttt{arXiv:2509.24527}}
\author{Gonzalo Fuentes}
\institute{Magíster en Ciencias de la Computación}
\date{\today}

% ─────────────────────────────────────────────────────────────────────────────
\begin{document}

% ── Portada ──────────────────────────────────────────────────────────────────
\begin{frame}
\titlepage
\end{frame}

% ── Índice ───────────────────────────────────────────────────────────────────
\begin{frame}{Contenidos}
\tableofcontents
\end{frame}

% ─────────────────────────────────────────────────────────────────────────────
\section{Motivación y Contexto}
% ─────────────────────────────────────────────────────────────────────────────

\begin{frame}{¿Por qué Modelos de Mundo?}
\begin{columns}[T]
  \begin{column}{0.5\textwidth}
    \textbf{El problema del RL model-free:}
    \begin{itemize}
      \item Requiere millones de interacciones reales con el entorno
      \item Muestras costosas en robots físicos o simulaciones complejas
      \item No aprovecha la estructura temporal de las observaciones
    \end{itemize}
    \vspace{8pt}
    \textbf{La promesa de los modelos de mundo:}
    \begin{itemize}
      \item Aprender una representación latente del entorno
      \item Predecir consecuencias de acciones \emph{sin ejecutarlas}
      \item Entrenar la política en \alert{imaginación} — infinitas muestras gratis
    \end{itemize}
  \end{column}
  \begin{column}{0.48\textwidth}
    \begin{block}{Línea de investigación}
      \centering
      \vspace{4pt}
      \textbf{DreamerV1} (2019) $\rightarrow$ aprendizaje por RSSM\\[4pt]
      \textbf{DreamerV2} (2020) $\rightarrow$ representaciones discretas\\[4pt]
      \textbf{DreamerV3} (2023) $\rightarrow$ generalización multi-tarea\\[4pt]
      \textbf{DreamerV4} (2025) $\rightarrow$ \alert{transformers + flow matching}
    \end{block}
    \vspace{8pt}
    \small DreamerV4 escala a imágenes de alta resolución y control de largo horizonte, superando a DreamerV3 y TD-MPC2.
  \end{column}
\end{columns}
\end{frame}

\begin{frame}{Desafíos del Control Visual a Largo Plazo}
\begin{columns}[T]
  \begin{column}{0.55\textwidth}
    \begin{enumerate}
      \item \textbf{Alta dimensionalidad:} imágenes $128{\times}128{\times}3$ son espacio de observación enorme
      \vspace{6pt}
      \item \textbf{Dependencias temporales largas:} las recompensas pueden estar distantes en el tiempo de las acciones que las causaron
      \vspace{6pt}
      \item \textbf{Eficiencia de muestras:} los entornos físicos/costosos requieren aprender con pocos datos reales
      \vspace{6pt}
      \item \textbf{Calidad de imaginación:} errores en el modelo se acumulan en rollouts de largo horizonte (\emph{compounding errors})
    \end{enumerate}
  \end{column}
  \begin{column}{0.42\textwidth}
    \begin{alertblock}{Contribución central}
      DreamerV4 propone un \textbf{pipeline de 4 fases} que separa explícitamente:
      \begin{itemize}
        \item Percepción (tokenizer MAE)
        \item Dinámica (flow matching)
        \item Comportamiento (finetune BC+reward)
        \item Política (PMPO en imaginación)
      \end{itemize}
    \end{alertblock}
  \end{column}
\end{columns}
\end{frame}

% ─────────────────────────────────────────────────────────────────────────────
\section{Visión General del Sistema}
% ─────────────────────────────────────────────────────────────────────────────

\begin{frame}{Pipeline de 4 Fases}
\centering
\begin{tikzpicture}[
  box/.style={rectangle, rounded corners=6pt, minimum width=2.8cm, minimum height=1.1cm,
              text centered, font=\small\bfseries, text=white},
  arr/.style={-{Latex[length=4pt]}, thick},
  node distance=0.55cm
]
  \node[box, fill=phaseA] (tok)  {\faEye\\ Tokenizer};
  \node[box, fill=phaseB, right=of tok] (dyn) {\faBrain\\ Dynamics};
  \node[box, fill=phaseC, right=of dyn] (ft)  {\faGraduationCap\\ Finetune};
  \node[box, fill=phaseD, right=of ft]  (agt) {\faRobot\\ Agent};

  \draw[arr] (tok) -- (dyn);
  \draw[arr] (dyn) -- (ft);
  \draw[arr] (ft)  -- (agt);

  \node[above=0.3cm of tok, font=\footnotesize, text=phaseA] {Fase 1a};
  \node[above=0.3cm of dyn, font=\footnotesize, text=phaseB] {Fase 1b};
  \node[above=0.3cm of ft,  font=\footnotesize, text=phaseC] {Fase 2};
  \node[above=0.3cm of agt, font=\footnotesize, text=phaseD] {Fase 3};

  \node[below=0.35cm of tok, font=\tiny, text=darkgray, align=center] {MAE\\compresión};
  \node[below=0.35cm of dyn, font=\tiny, text=darkgray, align=center] {Flow\\matching};
  \node[below=0.35cm of ft,  font=\tiny, text=darkgray, align=center] {BC +\\reward};
  \node[below=0.35cm of agt, font=\tiny, text=darkgray, align=center] {PMPO\\imaginación};

  % Ciclo de recolección
  \node[box, fill=darkgray, minimum width=12.8cm, minimum height=0.75cm,
        below=1.1cm of dyn] (collect) {\faDatabase\quad Recolección de datos (ciclo activo: política aleatoria $\to$ agente)};

  \draw[arr, gray] (agt.south) -- ++(0,-0.42) -- ++(-11.15,0) -- (collect.north west);
  \draw[arr, gray] (collect.north) -- ++(0,0.42);
\end{tikzpicture}

\vspace{10pt}
\begin{columns}[T]
  \begin{column}{0.24\textwidth}
    \small\textbf{\color{phaseA}Tokenizer:} codifica frames en tokens latentes comprimidos
  \end{column}
  \begin{column}{0.24\textwidth}
    \small\textbf{\color{phaseB}Dynamics:} predice estados futuros en espacio latente
  \end{column}
  \begin{column}{0.24\textwidth}
    \small\textbf{\color{phaseC}Finetune:} aprende comportamiento y valor desde datos reales
  \end{column}
  \begin{column}{0.24\textwidth}
    \small\textbf{\color{phaseD}Agent:} mejora política en rollouts imaginados
  \end{column}
\end{columns}
\end{frame}

% ─────────────────────────────────────────────────────────────────────────────
\section{Tokenizer (Fase 1a)}
% ─────────────────────────────────────────────────────────────────────────────

\begin{frame}{Tokenizer: Arquitectura MAE Temporal}
\begin{columns}[T]
  \begin{column}{0.54\textwidth}
    \textbf{Objetivo:} comprimir frames de observación en representaciones latentes compactas.

    \vspace{8pt}
    \textbf{Arquitectura:}
    \begin{itemize}
      \item \textbf{Patchificación:} imagen $H{\times}W$ dividida en parches $p{\times}p$ píxeles
      \item \textbf{Encoder} (Transformer bloque-causal):
            \begin{itemize}\small
              \item Atención espacial: latentes ven todos los parches; parches ven solo su modalidad
              \item Atención temporal: causal (pasado + presente)
            \end{itemize}
      \item \textbf{Bottleneck:} proyección a $d_b$ dims + $\tanh$ — fuerza compresión
      \item \textbf{Decoder:} reconstruye parches desde latentes via cross-attention
    \end{itemize}
  \end{column}
  \begin{column}{0.44\textwidth}
    \begin{block}{Entrenamiento MAE}
      \textbf{Masked Autoencoder:}
      \begin{itemize}
        \item Enmascara 0--90\% de parches aleatoriamente
        \item El encoder opera \emph{solo} sobre parches visibles + tokens latentes
        \item El decoder reconstruye \textbf{todos} los parches
        \item Pérdida: MSE + LPIPS (perceptual)
      \end{itemize}
    \end{block}
    \vspace{6pt}
    \begin{exampleblock}{Ventaja del masking agresivo}
      Con 90\% enmascarado, el bottleneck \emph{debe} comprimir información global del frame — imposible hacer trampa viendo los parches locales
    \end{exampleblock}
  \end{column}
\end{columns}
\end{frame}

\begin{frame}{Tokenizer: Compresión y Representación Espacial}
\begin{columns}[T]
  \begin{column}{0.5\textwidth}
    \textbf{Flujo de datos:}
    \vspace{4pt}

    \begin{tikzpicture}[scale=0.85, every node/.style={font=\scriptsize}]
      \node[draw, fill=phaseA!20, rounded corners, minimum width=3cm, minimum height=0.6cm]
           (img) at (0,0) {Frame $128{\times}128{\times}3$};
      \node[draw, fill=phaseA!20, rounded corners, minimum width=3cm, minimum height=0.6cm]
           (pat) at (0,-1.0) {$N_p$ parches, $d_p$ dims};
      \node[draw, fill=phaseA!40, rounded corners, minimum width=3cm, minimum height=0.6cm]
           (enc) at (0,-2.0) {Encoder Transformer};
      \node[draw, fill=orange!30, rounded corners, minimum width=3cm, minimum height=0.6cm]
           (bot) at (0,-3.0) {$N_L \times d_b$ latentes};
      \node[draw, fill=phaseB!30, rounded corners, minimum width=3cm, minimum height=0.6cm]
           (pak) at (0,-4.0) {Packing $\to N_z \times d_z$ espacial};

      \draw[-{Latex}, thick] (img)--(pat) node[midway,right,xshift=2pt]{patchify};
      \draw[-{Latex}, thick] (pat)--(enc) node[midway,right,xshift=2pt]{proyección};
      \draw[-{Latex}, thick] (enc)--(bot) node[midway,right,xshift=2pt]{bottleneck $+\tanh$};
      \draw[-{Latex}, thick] (bot)--(pak) node[midway,right,xshift=2pt]{reshape};
    \end{tikzpicture}
  \end{column}
  \begin{column}{0.48\textwidth}
    \textbf{Operación de Packing:}
    \vspace{6pt}

    Los latentes $N_L \times d_b$ se \emph{reempaquetan} en tokens espaciales para el modelo de dinámica:
    \[
      (B,\,T,\,N_L,\,d_b) \;\xrightarrow{\text{reshape}}\; (B,\,T,\,N_z,\,d_z)
    \]
    donde $N_z = N_L / k$ y $d_z = d_b \cdot k$

    \vspace{8pt}
    \begin{block}{Nuestra configuración}
      \begin{tabular}{ll}
        $p$ (patch size) & 4 px \\
        $N_p$ (parches)  & 1024 \\
        $N_L$ (latentes) & 32 \\
        $d_b$ (bottleneck) & 32 \\
        $k$ (packing)    & 2 \\
        $N_z$ (espacial) & 16 \\
        Compresión       & \textbf{32:1} \\
      \end{tabular}
    \end{block}
  \end{column}
\end{columns}
\end{frame}

% ─────────────────────────────────────────────────────────────────────────────
\section{Dynamics (Fase 1b)}
% ─────────────────────────────────────────────────────────────────────────────

\begin{frame}{Dynamics: Flow Matching con Shortcut Forcing}
\begin{columns}[T]
  \begin{column}{0.5\textwidth}
    \textbf{Objetivo:} dado el historial de estados latentes $z_{1:t}$ y acciones $a_{1:t}$, predecir $z_{t+1}$.

    \vspace{8pt}
    \textbf{¿Por qué Flow Matching?}
    \begin{itemize}
      \item Aprende distribuciones multimodales (el futuro puede ser ambiguo)
      \item Más estable que difusión: predice velocidad, no ruido
      \item Permite \emph{shortcuts}: inferencia en 1--4 pasos en lugar de 1000
    \end{itemize}

    \vspace{8pt}
    \textbf{Proceso de corrupción:}
    \[
      \tilde{z} = (1-\sigma)\,\epsilon + \sigma\, z^* , \quad \epsilon \sim \mathcal{N}(0,I)
    \]
    El modelo predice la \textbf{velocidad} $v = \frac{z^* - \tilde{z}}{1-\sigma}$
  \end{column}
  \begin{column}{0.48\textwidth}
    \begin{alertblock}{Shortcut Forcing}
      Innovación clave: impone consistencia \emph{entre niveles de ruido}.

      \vspace{4pt}
      Si el modelo predice bien $z^*$ desde $\tilde{z}_{\sigma}$, también debe predecir bien desde el punto intermedio $\tilde{z}_{\sigma/2}$.

      \vspace{4pt}
      Esto permite hacer inferencia con muy pocos pasos ODE en deployment ($k{=}2$--$4$ pasos vs $1000$ en DDPM).
    \end{alertblock}
    \vspace{6pt}
    \begin{block}{Arquitectura Dynamics}
      Transformer bloque-causal con:
      \begin{itemize}\small
        \item Tokens: \texttt{[ACT, SIGNAL, STEP, SPATIAL$\times 16$, REG$\times 4$]}
        \item Atención espacio-temporal separada
        \item $d_\text{model}{=}512$, profundidad $8$
      \end{itemize}
    \end{block}
  \end{column}
\end{columns}
\end{frame}

\begin{frame}{Dynamics: Entrenamiento y Bootstrapping}
\begin{columns}[T]
  \begin{column}{0.52\textwidth}
    \textbf{Función de pérdida:}
    \[
      \mathcal{L} = \underbrace{\mathcal{L}_\text{emp}}_{\text{datos reales}} + \lambda \underbrace{\mathcal{L}_\text{self}}_{\text{bootstrap}}
    \]

    \textbf{Pérdida empírica} ($\mathcal{L}_\text{emp}$):
    \begin{itemize}\small
      \item MSE entre velocidad predicha y verdadera
      \item Ponderada por $\sigma$ (mayor peso a pasos finos)
    \end{itemize}

    \vspace{4pt}
    \textbf{Pérdida de bootstrap} ($\mathcal{L}_\text{self}$):
    \begin{itemize}\small
      \item Activa después de calentar el modelo (paso $>$ bootstrap\_start)
      \item Usa las propias predicciones del modelo como objetivos
      \item Permite aprender transiciones no vistas en los datos
    \end{itemize}
  \end{column}
  \begin{column}{0.46\textwidth}
    \begin{exampleblock}{¿Para qué sirve el bootstrap?}
      \textbf{Sin bootstrap:} el modelo solo aprende transiciones observadas en los datos reales.

      \vspace{4pt}
      \textbf{Con bootstrap:} el modelo puede razonar sobre estados intermedios no observados, mejorando la calidad de los rollouts imaginados en el agente.
    \end{exampleblock}

    \vspace{8pt}
    \begin{block}{Tokens de registro (registers)}
      4 tokens aprendidos sin posición fija — actúan como ``scratch space'' para el transformer, mejorando representaciones globales.
    \end{block}
  \end{column}
\end{columns}
\end{frame}

% ─────────────────────────────────────────────────────────────────────────────
\section{Finetune y Agent (Fases 2 y 3)}
% ─────────────────────────────────────────────────────────────────────────────

\begin{frame}{Finetune: Comportamiento desde Datos Reales (Fase 2)}
\begin{columns}[T]
  \begin{column}{0.52\textwidth}
    \textbf{Objetivo:} adaptar el modelo de dinámica para predecir tanto el siguiente estado como la \emph{recompensa} y el \emph{valor esperado}.

    \vspace{8pt}
    \textbf{Pérdida combinada:}
    \[
      \mathcal{L}_\text{ft} = \underbrace{\mathcal{L}_\text{BC}}_{\text{imitar acciones}} + \lambda_r\underbrace{\mathcal{L}_\text{reward}}_{\text{predecir recompensa}}
    \]

    \vspace{4pt}
    \textbf{Behavioral Cloning} ($\mathcal{L}_\text{BC}$):
    \begin{itemize}\small
      \item Aprende a imitar la política de recolección de datos
      \item Resuelve el \emph{bootstrap problem}: da señal inicial al agente antes de que explore
    \end{itemize}

    \vspace{4pt}
    \textbf{Reward head:}
    \begin{itemize}\small
      \item Predice recompensa esperada desde el token de agente
      \item Necesario para que el agente pueda optimizar en imaginación
    \end{itemize}
  \end{column}
  \begin{column}{0.46\textwidth}
    \begin{alertblock}{El problema del bootstrap en RL}
      Una política aleatoria produce ventajas $\approx 0$ $\Rightarrow$ PMPO no puede mejorar.

      \vspace{4pt}
      BC ``ceba la bomba'': el agente parte sabiendo qué hacer razonablemente, y PMPO puede refinarlo.
    \end{alertblock}

    \vspace{8pt}
    \begin{block}{Arquitectura}
      Misma backbone que Dynamics + cabezas:
      \begin{itemize}\small
        \item \textbf{Policy head}: $\pi(a|z)$ — distribución de acciones
        \item \textbf{Value head}: $V(z)$ — retorno esperado
        \item \textbf{Reward head}: $r(z)$ — recompensa inmediata
      \end{itemize}
    \end{block}
  \end{column}
\end{columns}
\end{frame}

\begin{frame}{Agent: Entrenamiento por Imaginación con PMPO (Fase 3)}
\begin{columns}[T]
  \begin{column}{0.52\textwidth}
    \textbf{Policy Mirror Prox Optimization (PMPO):}
    \begin{itemize}
      \item Variante de TRPO adaptada para espacio de acciones continuas
      \item Actualización proximal: limita cuánto puede cambiar la política en cada paso
      \item Estable frente a ruido en las predicciones del modelo
    \end{itemize}

    \vspace{8pt}
    \textbf{Ciclo de imaginación:}
    \begin{enumerate}\small
      \item Muestrea estado inicial $z_0$ del replay buffer
      \item Despliega horizonte $H$ pasos usando Dynamics
      \item Calcula retornos con Value head + Reward head
      \item Actualiza política con PMPO
    \end{enumerate}
  \end{column}
  \begin{column}{0.46\textwidth}
    \begin{exampleblock}{Ventaja de la imaginación}
      Por cada interacción real con el entorno, el agente puede hacer \textbf{miles} de actualizaciones de política en imaginación — eficiencia de muestras extrema.
    \end{exampleblock}

    \vspace{8pt}
    \begin{block}{Pipeline activo (ciclos)}
      \centering
      \small
      Recolectar $\xrightarrow{}$ Tokenizer $\xrightarrow{}$ Dynamics $\xrightarrow{}$ Finetune $\xrightarrow{}$ Agent $\xrightarrow{}$ \textbf{Ciclo}
      \vspace{4pt}

      El agente mejora cada ciclo: datos mejores $\to$ modelo mejor $\to$ política mejor $\to$ datos mejores...
    \end{block}
  \end{column}
\end{columns}
\end{frame}

% ─────────────────────────────────────────────────────────────────────────────
\section{Innovaciones Técnicas Clave}
% ─────────────────────────────────────────────────────────────────────────────

\begin{frame}{Transformer Bloque-Causal}
\begin{columns}[T]
  \begin{column}{0.55\textwidth}
    \textbf{Problema:} los transformers estándar son cuadráticamente costosos en la secuencia completa.

    \vspace{8pt}
    \textbf{Solución:} arquitectura bloque-causal con dos tipos de atención:

    \vspace{6pt}
    \begin{tabular}{p{2.2cm}p{4.2cm}}
      \toprule
      \textbf{Tipo} & \textbf{Descripción} \\
      \midrule
      Espacial & Todos los tokens dentro del mismo timestep se ven entre sí (full attention) \\
      \addlinespace
      Temporal & Cada token solo ve tokens del mismo tipo en timesteps anteriores (causal) \\
      \bottomrule
    \end{tabular}

    \vspace{8pt}
    \small Los latentes ven todos los parches espacialmente; los parches solo ven su modalidad. Esto reduce el costo vs.\ atención global completa.
  \end{column}
  \begin{column}{0.43\textwidth}
    \begin{tikzpicture}[scale=0.7, font=\scriptsize]
      % Timesteps
      \foreach \t/\label in {0/$t{-}2$, 1/$t{-}1$, 2/$t$} {
        \node[draw, fill=phaseA!30, minimum width=1.5cm, minimum height=0.5cm, rounded corners]
          (lat\t) at (\t*2.1, 1.2) {Latentes};
        \node[draw, fill=gray!20, minimum width=1.5cm, minimum height=0.5cm, rounded corners]
          (pat\t) at (\t*2.1, 0.0) {Parches};
        \node[above=2pt of lat\t, font=\tiny] {\label};
      }
      % Spatial arrows (within timestep)
      \foreach \t in {0,1,2} {
        \draw[{Latex}-{Latex}, phaseA, thick] (lat\t) -- (pat\t);
      }
      % Temporal arrows (causal, latent to latent)
      \draw[-{Latex}, phaseB, thick] (lat0) to[bend left=25] (lat1);
      \draw[-{Latex}, phaseB, thick] (lat1) to[bend left=25] (lat2);
      \draw[-{Latex}, phaseB, thick] (lat0) to[bend left=40] (lat2);

      \node[right=0.2cm of lat2, text=phaseA, font=\scriptsize, align=left] {Espacial\\(bidirec.)};
      \node[right=0.2cm of pat2, text=phaseB, font=\scriptsize, align=left] {Temporal\\(causal)};
    \end{tikzpicture}

    \vspace{8pt}
    \begin{block}{Eficiencia}
      \small Profundidad $8$ capas, alternando atención espacial y temporal cada \texttt{time\_every}$=4$ capas.
    \end{block}
  \end{column}
\end{columns}
\end{frame}

\begin{frame}{Resumen de Innovaciones}
\begin{columns}[T]
  \begin{column}{0.48\textwidth}
    \begin{block}{\faEye\; Tokenizer MAE Temporal}
      \begin{itemize}\small
        \item Masking agresivo (hasta 90\%) fuerza representaciones comprimidas
        \item Bottleneck $\tanh$ previene colapso de representaciones
        \item Pérdida perceptual LPIPS para reconstrucciones de alta calidad
      \end{itemize}
    \end{block}
    \vspace{6pt}
    \begin{block}{\faBrain\; Flow Matching + Shortcuts}
      \begin{itemize}\small
        \item Velocidad en vez de ruido: más estable
        \item Shortcut forcing: inferencia en 2--4 pasos ODE
        \item Bootstrap self-supervised para estados no observados
      \end{itemize}
    \end{block}
  \end{column}
  \begin{column}{0.48\textwidth}
    \begin{block}{\faGraduationCap\; Finetune BC+Reward}
      \begin{itemize}\small
        \item Resuelve el cold-start del agente
        \item Aprende el modelo de recompensa necesario para PMPO
        \item Permite continual learning con datos cíclicos
      \end{itemize}
    \end{block}
    \vspace{6pt}
    \begin{block}{\faRobot\; PMPO en Imaginación}
      \begin{itemize}\small
        \item Miles de actualizaciones de política por interacción real
        \item Actualización proximal: estable frente al error del modelo
        \item Mejora continua en pipeline activo multi-ciclo
      \end{itemize}
    \end{block}
  \end{column}
\end{columns}
\end{frame}

% ─────────────────────────────────────────────────────────────────────────────
\section{Resultados Experimentales}
% ─────────────────────────────────────────────────────────────────────────────

\begin{frame}{Evaluación: DMControl y MMBench}
\begin{columns}[T]
  \begin{column}{0.52\textwidth}
    \textbf{Entornos de evaluación:}
    \begin{itemize}
      \item \textbf{DMControl}: locomotion continua (cheetah-run, walker-run, hopper-hop, etc.)
      \item \textbf{MMBench}: tareas con múltiples objetos y manipulación
      \item Hasta \textbf{30 tareas} simultáneas en el paper oficial
    \end{itemize}

    \vspace{8pt}
    \textbf{Baselines comparados:}
    \begin{itemize}
      \item \textbf{DreamerV3}: predecesor directo
      \item \textbf{TD-MPC2}: model-based state-of-the-art
      \item \textbf{SAC/PPO}: model-free references
    \end{itemize}
  \end{column}
  \begin{column}{0.46\textwidth}
    \begin{alertblock}{Resultados principales}
      \begin{itemize}
        \item DreamerV4 supera a DreamerV3 en \textbf{eficiencia de muestras}
        \item Mejor que TD-MPC2 en tareas de largo horizonte
        \item Escala favorablemente con más ciclos activos
        \item Los shortcuts reducen el tiempo de inferencia en ${\sim}10\times$
      \end{itemize}
    \end{alertblock}

    \vspace{6pt}
    \begin{exampleblock}{Ablación de $N_z$ (tokens espaciales)}
      \small $N_z{=}64$ da mejor balance — más tokens mejoran la predicción de dinámica pero aumentan el costo de atención cuadráticamente.
    \end{exampleblock}
  \end{column}
\end{columns}
\end{frame}

% ─────────────────────────────────────────────────────────────────────────────
\section{Nuestra Implementación}
% ─────────────────────────────────────────────────────────────────────────────

\begin{frame}{Implementación: Arquitectura Adoptada}
\begin{columns}[T]
  \begin{column}{0.5\textwidth}
    \textbf{Configuración base (128×128):}

    \vspace{6pt}
    \begin{tabular}{lcc}
      \toprule
      \textbf{Parámetro} & \textbf{Official} & \textbf{Nuestro} \\
      \midrule
      patch size        & 4        & \textbf{4} \\
      $N_L$ (latentes)  & 16       & \textbf{32} \\
      $d_b$ (bottleneck)& 32       & \textbf{32} \\
      $d_\text{model}$  & 256      & \textbf{256} \\
      depth (enc/dec)   & 8        & \textbf{8} \\
      dyn\_depth        & 8        & \textbf{8} \\
      $N_z$ (espacial)  & 8        & \textbf{16} \\
      Compresión        & 64:1     & \textbf{32:1} \\
      \bottomrule
    \end{tabular}
  \end{column}
  \begin{column}{0.48\textwidth}
    \textbf{Mejoras sobre el repo oficial:}
    \begin{itemize}
      \item \textbf{2× más latentes} ($N_L{=}32$ vs $16$): más capacidad representacional
      \item \textbf{Warmup + cosine LR}: convergencia más estable
      \item \textbf{Gradient clipping}: evita explosión de gradientes en Transformer
      \item \textbf{Pipeline activo}: el repo oficial no implementa finetune ni agente — nosotros sí
      \item \textbf{LR optimizado}: $3{\times}10^{-4}$ tokenizer, $10^{-4}$ dynamics
    \end{itemize}

    \vspace{6pt}
    \begin{block}{Hardware}
      2× NVIDIA H100 80GB\\
      \texttt{srun -{}-gpus=2 -{}-mem=128G}
    \end{block}
  \end{column}
\end{columns}
\end{frame}

\begin{frame}{Implementación: Pipeline Activo Multi-Ciclo}
\begin{columns}[T]
  \begin{column}{0.52\textwidth}
    \textbf{Tareas:} walker-run, cheetah-run, hopper-hop (DMControl)

    \vspace{6pt}
    \textbf{Por ciclo:}
    \begin{itemize}
      \item 24 episodios × 3 tareas × 1000 pasos = \textbf{72k frames}
      \item Ciclo 0: política aleatoria
      \item Ciclos 1+: política aprendida
    \end{itemize}

    \vspace{6pt}
    \textbf{Training cold (ciclo 0):}
    \begin{tabular}{lrr}
      \toprule
      Fase & Steps & Épocas \\
      \midrule
      Tokenizer  & 22\,500 & 10 \\
      Dynamics   & 15\,000 & 10 \\
      Finetune   & 15\,000 & 10 \\
      Agent PMPO & 10\,000 & --- \\
      \bottomrule
    \end{tabular}
  \end{column}
  \begin{column}{0.46\textwidth}
    \begin{exampleblock}{Resultados tokenizer (ciclo 0, 1 época)}
      Tras solo 750 pasos de entrenamiento ya se observan:
      \begin{itemize}\small
        \item Reconstrucciones coherentes del agente
        \item Pose corporal identificable
        \item Fondo (cielo/suelo) correctamente reconstruido
      \end{itemize}
    \end{exampleblock}

    \vspace{6pt}
    \begin{block}{Scheduler LR (nuevo)}
      \small
      Steps 0--2000: warmup lineal $0 \to 3{\times}10^{-4}$\\
      Steps 2000+: cosine decay hasta $\approx 0$
    \end{block}
  \end{column}
\end{columns}
\end{frame}

% ─────────────────────────────────────────────────────────────────────────────
\section{Conclusiones}
% ─────────────────────────────────────────────────────────────────────────────

\begin{frame}{Conclusiones y Trabajo Futuro}
\begin{columns}[T]
  \begin{column}{0.5\textwidth}
    \textbf{DreamerV4 representa:}
    \begin{itemize}
      \item La aplicación de arquitecturas modernas (ViT/Transformers) al aprendizaje por refuerzo basado en modelos del mundo
      \item Una separación limpia entre percepción, dinámica y comportamiento
      \item Flow matching como alternativa superior a la difusión clásica para modelado de estados futuros
      \item Un pipeline activo que mejora iterativamente con cada ciclo de datos
    \end{itemize}
  \end{column}
  \begin{column}{0.48\textwidth}
    \textbf{Trabajo futuro (nuestra implementación):}
    \begin{itemize}
      \item Completar los 20 ciclos activos de entrenamiento
      \item Evaluar mejora en score DMControl por ciclo
      \item Explorar más tareas (manipulación, navegación)
      \item Comparar con DreamerV3 y TD-MPC2 en condiciones controladas
    \end{itemize}

    \vspace{8pt}
    \begin{alertblock}{Estado actual}
      Tokenizer funcionando — reconstrucciones de alta calidad a 128×128 px después de solo 1 época de entrenamiento.
    \end{alertblock}
  \end{column}
\end{columns}
\end{frame}

\begin{frame}[plain]
  \centering
  \vspace{2cm}
  {\Huge \textbf{¿Preguntas?}}\\[1.5cm]
  {\large Hansen et al. (2025)}\\[0.3cm]
  {\normalsize \texttt{arXiv:2509.24527}}\\[0.3cm]
  {\normalsize \texttt{github.com/nicklashansen/dreamer4}}\\[1.5cm]
  {\large Gonzalo Fuentes\\Magíster en Ciencias de la Computación}
\end{frame}

\end{document}
